\section{Related work}
\paragraph{Temporal graphs and event generation.}
GraphMixer illustrates strong temporal encoding with simple components
\citep{cong2023}; our Mixer-style survival predictor is an independent
adaptation, not a reproduction of its link-prediction protocol. CEP3 models
community event times \citep{wang2022}, and DG-Gen generates continuous-time
dynamic graphs \citep{hosseini2025}. Their generation targets differ from
conditional first exits of an already active population. We make no numerical
superiority claim over these systems. Such a comparison requires explicit
state initialization, unknown future-count stopping and outcome adapters.

\paragraph{Joint forecasting and reconciliation.}
TACTiS and TACTiS-2 learn flexible multivariate copulas
\citep{drouin2022,ashok2024}; separating learned marginals from spatial copulas
is also established \citep{huk2025}. Our positive-factor Gaussian and
Marshall--Olkin controls \citep{marshall1967} are deliberately restrictive
classical families. Typed current incidence specifies shared randomness for
changing entities; neither that construction nor the marginal--dependence
decomposition is claimed as a new statistical principle.

Probabilistic reconciliation already conditions forecast distributions on
aggregation constraints, including discrete counts \citep{zambon2024}.
We adapt a two-level importance-conditioning step to vectors of first-exit
category counts and compare it with direct Dirichlet--multinomial forecasts
and bottom-up aggregation. This is an established-method comparator, not a
new reconciliation algorithm or an exact reproduction of an entire software
package. Proper multivariate scores complement marginal scores; the limited
dependence sensitivity of energy scores motivates an additional variogram
score \citep{scheuerer2015}.

\paragraph{Operational forecasting.}
RIMS combines learned predictions and intercase simulation
\citep{meneghello2025}, while AgentSimulator models resource interactions
\citep{kirchdorfer2025}. Closely related work initializes business-process
simulation from current states and estimates conditional uncertainty,
including on BPI 2017 \citep{avramenko2026}. Current-state forecasting and
operational uncertainty are therefore established problems. Our empirical
distinction is to hold individual probabilities fixed when testing incidence,
then separately test correction of those probabilities, direct count
alternatives and joint group-risk scores across workflow and compute logs.
Executable workload generation, such as Mimesys \citep{kim2026}, has a
different output contract: our first-exit scenarios do not reproduce full
workloads, future arrivals or repeated process transitions.
